\section{Almost straight coordinates and null trace pairs}
\label{sec:conductivity}

Fix a compact totally disconnected set $B\subset\C$, put
$G=\Chat\setminus B$, and fix a finite list of tests satisfying
\eqref{eq:pair-null}. Let $\mathcal M$ be its nonempty class of models
from Definition~\ref{def:model}. Write
\[
 h(z)=z+\frac{b(h)}z+O(z^{-2}),\qquad
 \beta=\sup_{h\in\mathcal M}\Re b(h).
\]
Lemma~\ref{lem:normalization} makes $\beta$ finite and places all model
omitted sets in one fixed disk. In this section $w=X+iY$ denotes the
target coordinate.

We construct a height coordinate close to the usual vertical coordinate
on compact sets, with a small global gradient error.
The subsequent level-set argument will allow it to be combined with the
retained tests while preserving the pairwise trace condition.

\begin{proposition}\label{prop:height}
Suppose $d>0$ and $h\in\mathcal M$ satisfies $\Re b(h)>\beta-d$.
Set $F=\C\setminus h(G)$. There is a continuous
$v\in W^{1,2}_{\loc}(\C)$, constant on every component of $F$, such that
$\nabla v=0$ almost everywhere on $F$ and
\begin{equation}\label{eq:height-energy}
 \int_{\C}|\nabla(v-Y)|^2\,\dd A\le 2\pi d.
\end{equation}
The function $v$ is a locally uniform limit of imaginary coordinates of
normalized plane homeomorphisms conformal outside a fixed disk.
Moreover, for each compact $K\subset\C$, the quantities
$\|v-Y\|_{C(K)}$ tend to zero as $d\downarrow0$, uniformly over these
choices of nearly maximizing models. The exterior bounds used here
depend only on $B$.
\end{proposition}

\begin{proof}
For an integer $n\ge2$, set $k=n$ on $F$ and $k=1$ elsewhere. We first
construct the principal quasiconformal coordinate $\phi_n=a_n+iv_n$
satisfying
\begin{equation}\label{eq:conductivity-system}
 (a_n)_X=k(v_n)_Y,\qquad (a_n)_Y=-k(v_n)_X,
 \qquad
 \phi_n(w)=w+\frac{c_n}{w}+O(w^{-2}).
\end{equation}
Indeed this system is equivalent to the real-linear Beltrami equation
\[
 (\phi_n)_{\bar w}=q\,\overline{(\phi_n)_w},
 \qquad q=\frac{k-1}{k+1}.
\]
We use the classical Cauchy--Beurling construction of the Beltrami
solution \cite[Sections~1--3]{AhlforsBers1960}.  Let $\mathcal C$ be the plane Cauchy transform and
$\mathcal S=\partial\mathcal C$ its $L^2$-isometric Beurling transform,
with $\bar\partial\mathcal C g=g$. The equation
\[
 g=q\,\overline{1+\mathcal Sg}
\]
is an affine contraction on $L^2(\C)$ considered as a real Banach space.
Its solution is supported on $F$. Thus $\phi_n=w+\mathcal Cg$ is locally
Sobolev and has the expansion in \eqref{eq:conductivity-system}; compact
support also makes $g$ integrable. Set
$\mu=(\phi_n)_{\bar w}/(\phi_n)_w$ where the denominator is nonzero,
and $\mu=0$ otherwise. The distortion inequality
$|(\phi_n)_{\bar w}|\le\|q\|_\infty|(\phi_n)_w|$ shows that
$(\phi_n)_{\bar w}=\mu(\phi_n)_w$ almost everywhere,
$\|\mu\|_\infty<1$, and $\mu=0$ off $F$. Standard planar weak-solution
regularity, the Measurable Riemann Mapping Theorem, and holomorphic
factorization \cite{AhlforsBers1960,Ahlfors2006,AstalaIwaniecMartin2009} give a continuous representative
$\phi_n=H\circ\psi$, where $\psi$ is the principal quasiconformal
plane homeomorphism solving $\psi_{\bar w}=\mu\psi_w$ and $H$ is entire.
Since $\psi(w)=w+O(1/w)$ and $\phi_n$ has the displayed principal
expansion, $H(w)=w+O(1/w)$ near infinity. Liouville's Theorem makes
$H$ the identity. Consequently $\phi_n$ itself is a homeomorphism.
Direct differentiation
gives the signs and the Jacobian:
\begin{equation}\label{eq:conductivity-derivatives}
 (\phi_n)_w=\frac{k+1}{2}\bigl((v_n)_Y+i(v_n)_X\bigr),\quad
 (\phi_n)_{\bar w}=\frac{k-1}{2}\bigl((v_n)_Y-i(v_n)_X\bigr),\quad
 J_{\phi_n}=k|\nabla v_n|^2.
\end{equation}

The map $\phi_n\circ h$ is again a model. It is conformal on $G$, and
its new coefficient is $b(h)+c_n$. A homeomorphism carries complementary
components to complementary components, preserving points and Jordan
disks; uniform continuity on a compact set containing $F$ preserves
the null-family property. For an old continued test $U$, its new
continuation is exactly $U\circ\phi_n^{-1}$, with the same prescribed
constant on the corresponding component. Quasiconformal composition
preserves its global $W^{1,2}(\Chat)$ membership and its zero gradient
on the transformed complement. To see also the null-set issue in this
statement, approximate the continuous $U$ uniformly and strongly in
$W^{1,2}$ by smooth sphere functions. The quasiconformal Dirichlet
inequality makes their compositions converge in gradient norm, and
uniform convergence identifies the limit. The almost-everywhere chain
rule then applies. The forward Lusin N property of $\phi_n$ maps the
source-null exceptions to target-null sets, so the resulting gradient
vanishes almost everywhere on $\phi_n(F)$.
Approximate maximality therefore gives
\begin{equation}\label{eq:coefficient-gain}
 \Re c_n\le d.
\end{equation}

For $D$ larger than a disk containing $F$, the area formula and the
boundary Laurent expansion give
\[
 \int_{D_D}J_{\phi_n}\,\dd A
   =\area\bigl(\phi_n(D_D)\bigr)=\pi D^2+o(1)
       \quad(D\to\infty).
\]
Here $D_D=\{w:|w|<D\}$. The Divergence Theorem, or the weak boundary
formula on this smooth exterior circle, gives the exact identity
\[
 \int_{D_D}(v_n)_Y\,\dd A
   =\int_{\partial D_D}v_n\,\nu_Y\,\dd s
   =\pi D^2-\pi\Re c_n;
\]
all other Laurent terms have zero contribution. Expanding the square,
using \eqref{eq:conductivity-derivatives}, and then letting $D\to\infty$
yields
\begin{equation}\label{eq:conductivity-energy}
 \int_{\C}|\nabla(v_n-Y)|^2\,\dd A
 +(n-1)\int_F|\nabla v_n|^2\,\dd A
   =2\pi\Re c_n\le2\pi d.
\end{equation}
In particular both terms on the left are controlled, with the stated
normalization constant.

We next obtain compactness without a uniform quasiconformal distortion
bound. A coordinate of a plane homeomorphism attains its maximum and
minimum on the boundary of every relatively compact open set: an
interior coordinate extremum contradicts openness. For almost every
circle centered at $x$, the Sobolev circle restriction and
Cauchy--Schwarz therefore give, for $0<r<R$,
\begin{equation}\label{eq:coordinate-modulus}
 \bigl(\osc_{\overline{D(x,r)}}v_n\bigr)^2\log(R/r)
 \le 2\pi\int_{D(x,R)}|\nabla v_n|^2\,\dd A.
\end{equation}
One obtains this inequality by bounding the oscillation on the smaller
disk by that on a circle of radius $t\in(r,R)$, bounding the latter
by the integral of $|\nabla v_n|$ around that circle, and integrating
with respect to $\dd t/t$.
Equation~\eqref{eq:conductivity-energy} bounds the right-hand side
locally, uniformly in $n$. Exterior univalence and principal
normalization, on a common exterior disk, give
\begin{equation}\label{eq:height-exterior}
 |v_n(w)-Y|\le C/|w|\qquad(|w|\ge R_0),
\end{equation}
where $C,R_0$ depend only on the common bound for $F$, hence only on
$B$. The maximum and minimum principles give uniform bounds on compact
sets as well. Thus \eqref{eq:coordinate-modulus} gives local
equicontinuity.

Pass to a locally uniform and locally Sobolev-weak subsequential limit
$v$. Lower semicontinuity proves \eqref{eq:height-energy}, and the
second term of \eqref{eq:conductivity-energy} forces $\nabla v=0$
almost everywhere on $F$. The function is therefore constant on the
connected interior of each Jordan disk component, and continuity
extends the same value to its boundary. Singleton components need no
argument. The coordinate maximum and minimum principles and the
exterior estimate pass to the limit.

Finally, allow the nearly maximizing models to vary and let $d\to0$.
The same modulus and exterior bounds give compactness on every compact
set. Any subsequential limit of $v-Y$ has zero distributional gradient
by \eqref{eq:height-energy}, hence is constant. Estimate
\eqref{eq:height-exterior}, with arbitrarily large $|w|$, makes the
constant zero. This proves the uniform compact-set closeness asserted
in the proposition.
\end{proof}

\begin{lemma}\label{lem:level-pairs}
Let $F$ and $v$ be as in Proposition~\ref{prop:height}. Every point
$x\in\C$ lies in a compact nondegenerate continuum of its $v$-level
contained in any prescribed closed ball of positive radius centered at
$x$ and meeting that ball's boundary. If $U\in C(\C)\cap W^{1,2}_{\loc}(\C)$ is real-valued and
$\nabla U=0$ almost everywhere on $F$, then
\begin{equation}\label{eq:level-null-pairs}
 \area\bigl((U,v)(F)\bigr)=0,
 \qquad
 \area\bigl((U,X)(F)\bigr)=0.
\end{equation}
\end{lemma}

\begin{proof}
For the approximating homeomorphism $\phi_n$, the inverse image of the
horizontal line through $\phi_n(x)$ is a proper simple curve through
$x$. Stop one branch at its first exit from the prescribed ball.
The resulting compact arcs have a Hausdorff convergent subsequence.
Their limit is connected, contains $x$, and meets the ball boundary.
Local uniform convergence $v_n\to v$ places it in the actual level
$v^{-1}(v(x))$. This proves the first assertion pointwise.

We use the scalar Sobolev Coarea Theorem, due to Federer
\cite{Federer1969}, in the precise-representative form of
Mal\'y, Swanson, and Ziemer \cite[Theorem~1.1]{MSZ2003}. With domain dimension $2$, target
dimension $1$, and exponent $2$, it gives, for every nonnegative Borel
weight $a$ supported in a bounded region,
\begin{equation}\label{eq:scalar-coarea}
 \int_{\R}\int_{v^{-1}(t)}a\,\dd\mathcal H^1\,\dd t
   =\int_{\C}a|\nabla v|\,\dd A.
\end{equation}
Continuity makes $v$ equal everywhere to its precise representative,
so these are the full pointwise level sets. In particular, almost every
level has finite length in any fixed compact box. Formula
\eqref{eq:scalar-coarea} also implies that planar null sets have zero
length on almost every level.

Choose a closed box $Q$ with $F$ in its interior. On a neighborhood of
$Q$, smooth approximation of the continuous Sobolev function $U$
provides $U_j\to U$ uniformly on $Q$ with
\[
 \sum_j\|\nabla U_j-\nabla U\|_{L^2(Q)}<\infty.
\]
Let $g$ be a nonnegative Borel representative of $|\nabla U|$.
By \eqref{eq:scalar-coarea} and Cauchy--Schwarz, outside one null set
of values $t$ we have simultaneously
\begin{gather*}
 \mathcal H^1(Q\cap v^{-1}(t))<\infty,\qquad
 \int_{Q\cap v^{-1}(t)}g\,\dd\mathcal H^1<\infty,\\
 \sum_j\int_{Q\cap v^{-1}(t)}
       |\nabla U_j-\nabla U|\,\dd\mathcal H^1<\infty,
 \qquad
 \int_{F\cap v^{-1}(t)}g\,\dd\mathcal H^1=0.
\end{gather*}
The last assertion uses $g=0$ almost everywhere on $F$; the coarea
identity makes the choice on the planar-null exceptional set
irrelevant. For every such good level and \emph{every} rectifiable
simple arc $\gamma\subset Q\cap v^{-1}(t)$, the smooth arc inequality
passes to the limit to give
\begin{equation}\label{eq:level-arc}
 |U(x)-U(y)|\le\int_\gamma g\,\dd\mathcal H^1
 \qquad(x,y\text{ the endpoints of }\gamma).
\end{equation}
Indeed the gradient error on any arc is bounded by the error on the
entire level section, whereas uniform convergence controls both
endpoints. Thus the exceptional levels do not depend on the choice of
arc or its endpoints.

Fix a good level. Its nondegenerate components in $Q$ form an at most
countable family of finite-length continua. Each such continuum is
locally connected: failure of local connectedness, by boundary
bumping, would give infinitely many disjoint subcontinua of diameter
bounded below, contrary to finite length. A locally connected metric
continuum is locally arcwise connected, so its connected relatively
open subsets are arcwise connected. All their simple arcs are
rectifiable, since their lengths are bounded by the length of the
whole continuum.

Let $K$ be one of these components and $A=F\cap K$. The finite measure
$g\mathcal H^1|_K$ vanishes on $A$. Given $\epsilon>0$, choose a
relatively open neighborhood $O$ of $A$ in $K$ with
$\int_Og\,\dd\mathcal H^1<\epsilon$. Local connectedness and second
countability make its connected components $O_i$ a countable family
of relatively open sets. Joining any two points of $O_i$ by a simple
arc there and using \eqref{eq:level-arc} shows
\[
 \diam U(O_i)\le\int_{O_i}g\,\dd\mathcal H^1.
\]
Since each $U(O_i)$ is an interval, the outer length of $U(A)$ is at
most $\epsilon$, and hence is zero. Every point of
$F\cap v^{-1}(t)$ belongs to a nondegenerate component of
$Q\cap v^{-1}(t)$ by the first assertion of the lemma, using a small
ball inside $Q$. Thus the countable family of components covers
\emph{every} such point, and $U(F\cap v^{-1}(t))$ has length zero.
The compact set $(U,v)(F)$ now has almost every horizontal section of
length zero. Fubini proves its area is zero.

For the companion pair, the continuous Sobolev representative has the
absolute continuity on lines (ACL) property. Thus $U$ is absolutely
continuous on almost every
vertical line. Fubini gives $\partial_YU=0$ almost everywhere on that
line's intersection with $F$. The one-dimensional area inequality for
an absolutely continuous function therefore makes the image of this
intersection length-null. Applying Fubini to the compact set
$(U,X)(F)$ proves the second equality in
\eqref{eq:level-null-pairs}.
\end{proof}
